\documentclass[12pt,a4paper]{article}

\usepackage[utf8]{inputenc}
\usepackage[romanian]{babel}
\usepackage{geometry}
\geometry{margin=2.5cm}
\usepackage{graphicx}
\usepackage{booktabs}
\usepackage{hyperref}
\usepackage{listings}
\usepackage{xcolor}
\usepackage{amsmath}
\usepackage{pgfplotstable}
\usepackage{pgfplots}
\pgfplotsset{compat=1.18}
\usepackage{csvsimple}
\usepackage{float}

\lstset{
    language=Python,
    basicstyle=\ttfamily\small,
    keywordstyle=\color{blue},
    stringstyle=\color{red},
    commentstyle=\color{gray},
    numbers=left,
    numberstyle=\tiny\color{gray},
    breaklines=true,
    frame=single,
}

\title{
    \textbf{Travel Companion -- AI Restaurant Finder}\\[0.5em]
    \large Varianta Secvențială (S6)\\[0.3em]
    \large Algoritmi Paraleli și Distribuiți
}
\author{
    Stoian Cristian Georgian\\
    Calculatoare Română, Grupa CR3.2A
}
\date{Martie 2026}

\begin{document}

\maketitle

\tableofcontents
\newpage

% ============================================================
\section{Cerințe și Tema Proiectului}

\subsection{Descrierea Proiectului}

\textbf{Travel Companion} este o aplicație bazată pe inteligență artificială care recomandă cele mai bune restaurante pe baza preferințelor utilizatorului (locație, tip de mâncare, mediu, distanță etc.).

Sistemul agregă date de la multiple surse API externe și aplică algoritmi de scoring cu ponderi configurabile pentru a determina cele mai potrivite restaurante.

\subsection{Obiectivul Principal}

Obiectivul principal al proiectului este \textbf{analiza și compararea abordărilor secvențiale și paralele} pentru agregarea și procesarea datelor din surse multiple. Această livrare (S6) prezintă implementarea și analiza variantei secvențiale.

\subsection{Modele de Execuție}

\begin{enumerate}
    \item \textbf{Secvențial} (S6 -- această livrare): Toate operațiile se execută una după alta
    \item \textbf{Paralel \#1 -- Async I/O}: Apeluri API concurente folosind \texttt{asyncio}
    \item \textbf{Paralel \#2 -- Multiprocessing}: Procesare CPU paralelă pentru scoring
    \item \textbf{Paralel \#3 -- Hibrid}: Async I/O + Multiprocessing combinat
\end{enumerate}

\subsection{Tehnologii Utilizate}

\begin{itemize}
    \item \textbf{Limbaj}: Python 3.12.3
    \item \textbf{Pattern arhitectural}: Strategy Pattern (interfață comună pentru toate modurile de execuție)
    \item \textbf{Surse de date}: Clienți API reali (Geoapify, Overpass/OSM, Foursquare)
    \item \textbf{Instrumentare}: \texttt{TimingCollector} cu decorator \texttt{@timed\_task} pentru măsurarea timpilor
    \item \textbf{Dependențe}: \texttt{requests}, \texttt{python-dotenv}
\end{itemize}

\subsection{Surse de Date API}
\label{sec:api}

Sistemul utilizează trei surse de date reale pentru restaurante:

\begin{enumerate}
    \item \textbf{Geoapify Places API} -- oferă date despre restaurante incluzând rating (bazat pe popularitate), locație, categorii, nivel de preț și adresă formatată. Suportă filtrarea pe categorii specifice (e.g. \texttt{catering.restaurant.seafood}) și filtre geografice circulare.

    \item \textbf{Overpass API (OpenStreetMap)} -- furnizează date deschise din baza de date OpenStreetMap. Permite interogări complexe prin limbajul Overpass QL pentru restaurante cu filtrare pe tipul de bucătărie (\texttt{cuisine}). Extrage tag-uri de mediu (terasă, loc de joacă, seating exterior) și adrese din componentele OSM. Nu necesită cheie API.

    \item \textbf{Foursquare Places API} -- adaugă o a treia perspectivă cu căutare text-based și categorii detaliate. Autentificarea se face prin Bearer token.
\end{enumerate}

Toți clienții implementează aceeași interfață abstractă \texttt{BaseAPIClient}, permițând adăugarea sau înlocuirea transparentă a oricărei surse de date.

\subsection{Geocodare}

Înainte de interogarea surselor de restaurante, sistemul trebuie să determine coordonatele GPS ale orașului specificat de utilizator. Rezolvarea coordonatelor se face ierarhic:

\begin{enumerate}
    \item \textbf{Coordonate explicite} -- dacă utilizatorul furnizează \texttt{latitude}/\texttt{longitude} în query
    \item \textbf{Cache local} -- un dicționar predefinit cu coordonatele celor mai comune orașe românești (București, Constanța, Cluj-Napoca, Timișoara, Iași, Brașov)
    \item \textbf{Geoapify Geocoding API} -- pentru orice alt oraș, se apelează \texttt{/v1/geocode/search} cu parametrul \texttt{type=city}, permițând căutarea oricărui oraș din lume
    \item \textbf{Fallback} -- în cazul în care geocodarea eșuează, se folosesc coordonatele Bucureștiului ca valoare implicită
\end{enumerate}

Această abordare asigură funcționarea aplicației pentru orice locație, nu doar pentru orașele predefinite.

\subsection{Pipeline de Execuție}

Fluxul secvențial al pipeline-ului:

$$T_{total} = T_{geocode} + T_{API_1} + T_{API_2} + \ldots + T_{API_N} + T_{dedup} + T_{scoring} + T_{ranking}$$

Unde $T_{geocode}$ este timpul de rezolvare a coordonatelor (0 pentru orașele din cache, $\sim$0.3--0.5s pentru geocodare API), fiecare $T_{API_i}$ include latența de rețea și timpul de procesare a răspunsului, iar $T_{dedup}$, $T_{scoring}$, $T_{ranking}$ sunt timpii de procesare locală.

% ============================================================
\section{Arhitectura Sistemului}

\subsection{Strategy Pattern}

Arhitectura se bazează pe \textbf{Strategy Pattern}, unde interfața abstractă \texttt{ExecutionStrategy} definește metoda \texttt{execute()} implementată de fiecare mod de execuție:

\begin{lstlisting}[caption={Interfața ExecutionStrategy}]
class ExecutionStrategy(ABC):
    def __init__(self, collector: TimingCollector):
        self.collector = collector

    @abstractmethod
    def execute(self, query, clients, scorer):
        pass

    @property
    @abstractmethod
    def mode_name(self) -> str:
        pass
\end{lstlisting}

Pentru S6, este implementată doar \texttt{SequentialStrategy}. Variantele paralele se vor adăuga fără modificarea codului existent, datorită polimorfismului oferit de interfața comună.

\subsection{Structura Proiectului}

\begin{itemize}
    \item \texttt{models.py} -- Clase de date: \texttt{Query}, \texttt{Restaurant}, \texttt{ScoredRestaurant}, \texttt{PipelineResult}
    \item \texttt{config.py} -- Configurație: chei API, ponderi scoring, valori implicite
    \item \texttt{timing.py} -- \texttt{TimingCollector}, \texttt{TimingRecord} și decoratorul \texttt{@timed\_task}
    \item \texttt{api\_clients/} -- Clienți API:
    \begin{itemize}
        \item \texttt{base.py} -- Interfața abstractă \texttt{BaseAPIClient}
        \item \texttt{real\_client.py} -- Implementări reale: \texttt{GeoapifyClient}, \texttt{OverpassClient}, \texttt{FoursquareClient}
        \item \texttt{mock\_client.py} -- Clienți mock pentru teste (date sintetice cu latență simulată)
    \end{itemize}
    \item \texttt{execution/} -- Strategii de execuție:
    \begin{itemize}
        \item \texttt{base.py} -- Interfața abstractă \texttt{ExecutionStrategy}
        \item \texttt{sequential.py} -- Implementarea secvențială
    \end{itemize}
    \item \texttt{scoring.py} -- Formula de scoring cu ponderi configurabile și deduplicare
    \item \texttt{pipeline.py} -- Orchestratorul pipeline-ului
    \item \texttt{main.py} -- Punct de intrare CLI
    \item \texttt{experiments.py} -- Runner pentru experimente batch cu export CSV
\end{itemize}

\subsection{Modelul de Date}

\begin{lstlisting}[caption={Clasele principale de date}]
@dataclass
class Query:
    location: str           # "Constanta"
    food_types: List[str]   # ["seafood", "fish"]
    environment: List[str]  # ["seaside", "terrace"]
    max_distance_km: float  # 10.0
    latitude: Optional[float] = None
    longitude: Optional[float] = None

@dataclass
class Restaurant:
    name: str
    source: str             # "geoapify", "overpass_osm", "foursquare"
    rating: float           # 0-5
    latitude: float
    longitude: float
    distance_km: float
    categories: List[str]
    price_level: int        # 1-4
    environment_tags: List[str]
    review_count: int
    address: str
\end{lstlisting}

\subsection{Formula de Scoring}

$$score = 0.30 \cdot S_{rating} + 0.25 \cdot S_{distance} + 0.30 \cdot S_{food} + 0.15 \cdot S_{environment}$$

Unde:
\begin{itemize}
    \item $S_{rating} = \frac{rating}{5.0}$ -- normalizare la intervalul $[0, 1]$
    \item $S_{distance} = \max(0, 1 - \frac{d}{d_{max}})$ -- penalizare liniară cu distanța
    \item $S_{food}$ = raportul de potrivire fuzzy între tipurile de mâncare căutate și categoriile restaurantului (matching pe subșiruri)
    \item $S_{environment}$ = raportul de potrivire fuzzy între preferințele de mediu și tag-urile restaurantului
\end{itemize}

\textbf{Deduplicare:} Restaurantele cu același nume aflate la mai puțin de 100m distanță sunt considerate duplicate. Se păstrează versiunea cu rating-ul mai mare.

% ============================================================
\section{Specificații Mașină}

\begin{table}[H]
\centering
\begin{tabular}{ll}
\toprule
\textbf{Parametru} & \textbf{Valoare} \\
\midrule
CPU & AMD Ryzen 7 7745HX with Radeon Graphics \\
Nuclee fizice & 8 \\
Nuclee logice & 16 (8 cores $\times$ 2 threads) \\
RAM & 16 GB \\
Sistem de operare & Linux 6.6.87.2-microsoft-standard-WSL2 x86\_64 \\
Python & 3.12.3 \\
Conexiune rețea & Wi-Fi (apeluri API reale către servere externe) \\
\bottomrule
\end{tabular}
\caption{Specificațiile mașinii de test}
\label{tab:machine}
\end{table}

% ============================================================
\section{Implementare Secvențială}

\subsection{Descrierea Fluxului}

În modul secvențial, toate operațiile se execută strict una după alta:

\begin{enumerate}
    \item Se construiește query-ul din parametrii utilizatorului (locație, tip mâncare, mediu)
    \item Se rezolvă coordonatele GPS ale orașului (cache local sau geocodare API)
    \item Se apelează API \#1 (\textbf{Geoapify}) -- se așteaptă răspunsul complet
    \item Se apelează API \#2 (\textbf{Overpass/OSM}) -- se așteaptă răspunsul complet
    \item Se apelează API \#3 (\textbf{Foursquare}) -- se așteaptă răspunsul complet (opțional)
    \item Se deduplică rezultatele (eliminare duplicate pe bază de nume + locație)
    \item Se calculează scorurile pentru fiecare restaurant
    \item Se sortează restaurantele descrescător după scor
    \item Se returnează top K rezultate
\end{enumerate}

\begin{lstlisting}[caption={Implementarea SequentialStrategy}]
class SequentialStrategy(ExecutionStrategy):
    mode_name = "sequential"

    def execute(self, query, clients, scorer):
        all_restaurants = []
        for client in clients:
            start = time.perf_counter()
            try:
                results = client.search(query)
            except Exception as e:
                # Log error, continue with other APIs
                continue
            end = time.perf_counter()
            self.collector.record(
                f"api_{client.name}", start, end,
                count=len(results))
            all_restaurants.extend(results)

        # Deduplication, scoring, ranking
        deduped = scorer.deduplicate(all_restaurants)
        scored = scorer.score_all(deduped, query)
        ranked = sorted(scored,
                       key=lambda x: x.score,
                       reverse=True)
        return ranked
\end{lstlisting}

\subsection{Tratarea Erorilor}

Dacă un client API eșuează (timeout, eroare de rețea, răspuns invalid), execuția continuă cu ceilalți clienți. Eroarea este înregistrată în \texttt{TimingCollector} cu metadatele corespunzătoare, iar utilizatorul primește un avertisment.

\subsection{Instrumentare Timing}

Fiecare etapă este instrumentată cu \texttt{TimingCollector}, care înregistrează:
\begin{itemize}
    \item Timpul de start și end (relativ la începutul pipeline-ului), folosind \texttt{time.perf\_counter()}
    \item Durata în secunde
    \item Metadate: număr de rezultate returnate, mesaje de eroare (dacă există)
\end{itemize}

Datele de timing sunt exportate în format CSV prin metoda \texttt{to\_csv\_dict()}, care aplatizează înregistrările într-un dicționar cu coloane separate per etapă de pipeline.

% ============================================================
\section{Rezultate Experimentale}

\subsection{Configurarea Experimentelor}

Experimentele au fost rulate folosind \textbf{clienți API reali} (Geoapify, Overpass, Foursquare), cu apeluri de rețea efective către serverele externe.

\begin{itemize}
    \item \textbf{6 query-uri diferite} (orașe și tipuri de mâncare variate):
    \begin{itemize}
        \item Constanța / seafood / seaside
        \item București / italian + pizza / cozy
        \item Cluj-Napoca / romanian + traditional
        \item Timișoara / fast food / central
        \item Iași / french / elegant
        \item Brașov / steakhouse / rustic + terrace
    \end{itemize}
    \item 1, 2 și 3 surse API per configurație
    \item 3 rulări per query per configurație
    \item Total: $6 \times 3 \times 3 = 54$ rulări
\end{itemize}

\subsection{Rezultate -- Timpi de Execuție}

\begin{table}[H]
\centering
\begin{tabular}{lccc}
\toprule
\textbf{Nr. APIs} & \textbf{Media (s)} & \textbf{Std Dev (s)} & \textbf{Min / Max (s)} \\
\midrule
1 API (Geoapify)              & 0.6907 & 0.2442 & 0.4190 / 1.2432 \\
2 APIs (Geoapify + Overpass)  & 9.0550 & 3.5290 & 1.4212 / 14.1986 \\
3 APIs (+ Foursquare)         & 8.7294 & 4.5106 & 1.7785 / 16.7381 \\
\bottomrule
\end{tabular}
\caption{Timpi de execuție totali -- mod secvențial}
\label{tab:results}
\end{table}

\subsection{Timpi Medii per Sursă API}

\begin{table}[H]
\centering
\begin{tabular}{lcc}
\toprule
\textbf{Sursă API} & \textbf{Timp mediu (s)} & \textbf{Observații} \\
\midrule
Geoapify         & $\sim$0.69 -- 0.78 & Rapid, răspunsuri consistente \\
Overpass (OSM)   & $\sim$7.39 -- 8.35 & Foarte variabil, depinde de încărcarea serverului \\
Foursquare       & $\sim$0.56          & Rapid, răspunsuri consistente \\
\midrule
Deduplicare      & $<$ 0.001          & Neglijabil \\
Scoring          & $<$ 0.001          & Neglijabil \\
Ranking          & $<$ 0.001          & Neglijabil \\
\bottomrule
\end{tabular}
\caption{Timpi medii per componentă a pipeline-ului}
\label{tab:api_times}
\end{table}

\subsection{Analiza Rezultatelor}

\subsubsection{Dominanța latenței API}

Datele experimentale confirmă că \textbf{latența API domină complet timpul de execuție}. Timpii de procesare locală (deduplicare, scoring, ranking) sunt sub 1ms, reprezentând mai puțin de 0.01\% din timpul total.

$$T_{total} \approx \sum_{i=1}^{N} T_{API_i}$$

\subsubsection{Variabilitatea Overpass API}

API-ul Overpass (OpenStreetMap) prezintă cea mai mare variabilitate, cu timpi între 0.47s și 15.27s. Aceasta se datorează:
\begin{itemize}
    \item Încărcării serverelor publice Overpass (shared infrastructure)
    \item Complexității interogărilor Overpass QL (filtrare pe rază + tip bucătărie)
    \item Lipsei unui SLA garantat (serviciu gratuit, open-source)
\end{itemize}

\subsubsection{Scalabilitate Liniară}

Cu 1 API, timpul mediu este $\sim$0.69s. Cu 2 APIs, crește la $\sim$9.05s (dominat de Overpass). Cu 3 APIs, timpul rămâne similar ($\sim$8.73s) deoarece Foursquare este rapid ($\sim$0.56s) și se adaugă la un total deja dominat de Overpass.

Comportamentul este strict secvențial:
$$T_{2\text{ APIs}} \approx T_{\text{Geoapify}} + T_{\text{Overpass}} \approx 0.71 + 8.35 = 9.06s$$
$$T_{3\text{ APIs}} \approx T_{\text{Geoapify}} + T_{\text{Overpass}} + T_{\text{Foursquare}} \approx 0.78 + 7.39 + 0.56 = 8.73s$$

\subsubsection{Oportunitate de Paralelizare}

Apelurile API sunt \textbf{complet independente} -- nu există dependențe de date între ele. În modul secvențial, fiecare apel blochează execuția până la primirea răspunsului. Cu paralelizare (async I/O), timpul ar scădea la:

$$T_{paralel} \approx \max(T_{API_1}, T_{API_2}, T_{API_3})$$

Speedup-ul teoretic pentru 3 APIs:
$$S = \frac{T_{secvential}}{T_{paralel}} = \frac{T_1 + T_2 + T_3}{\max(T_1, T_2, T_3)} \approx \frac{8.73}{7.39} \approx 1.18\times$$

În cazul acesta, speedup-ul este modest deoarece Overpass domină. Pentru API-uri cu latențe echilibrate, speedup-ul ar fi mai aproape de $N$ (numărul de surse).

% ============================================================
\section{Utilizare}

\subsection{Interfața CLI}

\begin{lstlisting}[language=bash, caption={Exemple de utilizare}]
# Instalare dependente
pip install -r requirements.txt

# Cautare restaurante in Constanta
python main.py --location "Constanta" \
    --food "seafood,fish" \
    --environment "seaside,terrace" \
    --apis 2 --top-k 5

# Rulare experimente batch
python experiments.py --runs-per-query 5 \
    --api-counts "1,2,3" \
    --output report/experiment_results.csv
\end{lstlisting}

\subsection{Parametri CLI}

\begin{table}[H]
\centering
\begin{tabular}{lll}
\toprule
\textbf{Parametru} & \textbf{Implicit} & \textbf{Descriere} \\
\midrule
\texttt{--location} & (obligatoriu) & Numele orașului \\
\texttt{--food} & \texttt{""} & Tipuri de mâncare, separate prin virgulă \\
\texttt{--environment} & \texttt{""} & Preferințe de mediu, separate prin virgulă \\
\texttt{--max-distance} & 10.0 & Distanța maximă în km \\
\texttt{--apis} & 2 & Număr de surse API (1--3) \\
\texttt{--top-k} & 10 & Număr de rezultate afișate \\
\bottomrule
\end{tabular}
\caption{Parametrii liniei de comandă}
\label{tab:cli}
\end{table}

% ============================================================
\section{Concluzii și Direcții Viitoare}

\subsection{Concluzii S6}

\begin{itemize}
    \item Varianta secvențială funcționează corect cu API-uri reale și produce recomandări relevante
    \item Timpii de execuție sunt dominați de latența API (în special Overpass/OSM)
    \item Procesarea locală (deduplicare, scoring, ranking) este neglijabilă ($<$ 1ms)
    \item Variabilitatea mare a timpilor Overpass (0.47s -- 15.27s) subliniază importanța toleranței la erori
    \item Arhitectura Strategy Pattern facilitează adăugarea variantelor paralele fără modificarea codului existent
\end{itemize}

\subsection{Direcții Viitoare}

\begin{enumerate}
    \item \textbf{Paralel \#1 (Async I/O)}: Toate apelurile API lansate simultan cu \texttt{asyncio} + \texttt{aiohttp} -- $T = \max(T_{API_1}, T_{API_2}, T_{API_3})$
    \item \textbf{Paralel \#2 (Multiprocessing)}: Scoring distribuit pe multiple procese pentru seturi mari de date
    \item \textbf{Paralel \#3 (Hibrid)}: Combinarea async I/O (pentru API) + multiprocessing (pentru scoring)
    \item Vizualizare performanță: grafice Gantt pentru timeline-ul execuției, grafice speedup
\end{enumerate}

\end{document}
